
本小节余下部分证明命题~\ref{prop:7-2-275}、
\ref{prop:7-2-277}。
开始之前，先引入记号 $(n,K)$、$(n',K')$。
在两命题的证明中，
需要多次更换
$n'''(\lambda_i)$、$K'''(\lambda_i)$
及
$n''''(\lambda_i)$、$K''''(\lambda_i)$
（$i=0,1$）。
每次更换仍使
\[
\lim_{\lambda_i\to0}n'''(\lambda_i)
=\lim_{\lambda_i\to0}K'''(\lambda_i)
=\infty.
\]
还可选取它们，使
\[
\begin{aligned}
n'''(\lambda_i)&=n'''(\lambda'_i)
=n''''(\lambda_i)=n''''(\lambda'_i),\\
K'''(\lambda_i)&=K'''(\lambda'_i)
=K''''(\lambda_i)=K''''(\lambda'_i).
\end{aligned}
\]
但须注意，不能要求
$n'''(\lambda_i)=n'''(\lambda_{i+1})$。
这是因为证明定理~\ref{thm:4-6-25}
的策略，是利用命题~\ref{prop:7-2-275}、
\ref{prop:7-2-277}，
把 $A_{n'''(\lambda_i),K'''(\lambda_i)}$
结构扩张为
$A_{n'''(\lambda_{i+1}),K'''(\lambda_{i+1})}$
结构。

以下在本小节中简记
\[
\begin{aligned}
(n,K)
&=(n'''(\lambda_0),K'''(\lambda_0))
 =(n'''(\lambda'_0),K'''(\lambda'_0)),\\
(n',K')
&=(n'''(\lambda_1),K'''(\lambda_1))
 =(n'''(\lambda'_1),K'''(\lambda'_1)).
\end{aligned}
\]
证明过程中可以改变
$n'''(\lambda_0)$ 等的选取，
但不改变 $n$ 等记号。
回忆 $(n,K)<(n',K')$。
如果构造的结构只涉及
$L'_1$、$L_1$ 或
$[0,1]\times L'_1$ 的单形，
它属 $A_{n',K'}$ 类；
如果涉及 $L_0$、$L'_0$
或 $[0,1]\times L'_0$ 的单形，
它属 $A_{n,K}$ 类。

现在开始证明命题~\ref{prop:7-2-275}。
先回顾图式~\ref{diag:7-2-14} 中
同伦 $\mathfrak g_+^{\{\Phi_{r,s_0}\}_r}$
的构造。它是下列三个同伦的复合：
\begin{equation}
\mathfrak g^{\phi'}\circ
\mathfrak g^{\psi_{0,s_0}}
\sim\mathfrak g^{\phi'\circ\psi_{0,s_0}},
\tag{7.2.280.1}\label{eq:7-2-280-1}
\end{equation}
\begin{equation}
\mathfrak g^{\phi'\circ\psi_{0,s_0}}
\sim\mathfrak g^{\psi_{1,s_0}\circ\phi},
\tag{7.2.280.2}\label{eq:7-2-280-2}
\end{equation}
\begin{equation}
\mathfrak g^{\psi_{1,s_0}\circ\phi}
\sim\mathfrak g^{\psi_{1,s_0}}
\circ\mathfrak g^\phi.
\tag{7.2.280.3}\label{eq:7-2-280-3}
\end{equation}
其中 (7.2.280.1)、(7.2.280.3)
由定理~\ref{thm:4-6-44}
的 $A_{n,K}$ 版本定义；
(7.2.280.2) 由同伦
$\{\Phi_{r,s_0}\}_r$ 得到。
把 (7.2.280.2) 的同伦记为
$\mathfrak g^{\{\Phi_{r,s_0}\}_r}$。

我们要构造的同伦的同伦
$\mathfrak H$ 也由同伦的同伦胶合而成。
事实上，(7.2.280.1)、
(7.2.280.3) 各拆为两段，
因而 $\mathfrak H$ 由五段
同伦的同伦胶合得到。

回忆第 4.6 节先构造
$\mathfrak f^\phi$ 等，
再由此构造 $\mathfrak g^\phi$ 等。
现复看
$\mathfrak f^{\psi_{1,s_0}\circ\phi}$、
$\mathfrak f^{\phi'\circ\psi_{0,s_0}}$。
第 7.2.9 小节的论证给出滤过
$A_{n',K'}$ 代数
\[
C^{r=r_0,s=s_0}
(L'_1,(\lambda_0,\lambda'_1))
\]
及滤过 $A_{n,K}$ 同态
\begin{equation}
\mathfrak f^{\psi_{1,s_0}\circ\phi}\colon
C(L_0,\lambda_0)\longrightarrow
C^{r=1,s=s_0}(L'_1,(\lambda_0,\lambda'_1)),
\tag{7.2.281.1}\label{eq:7-2-281-1}
\end{equation}
\begin{equation}
\mathfrak f^{\phi'\circ\psi_{0,s_0}}\colon
C(L_0,\lambda_0)\longrightarrow
C^{r=0,s=s_0}(L'_1,(\lambda_0,\lambda'_1)).
\tag{7.2.281.2}\label{eq:7-2-281-2}
\end{equation}
于是
$\mathfrak g^{\psi_{1,s_0}\circ\phi}$、
$\mathfrak g^{\phi'\circ\psi_{0,s_0}}$
分别是 (7.2.281.1)、
(7.2.281.2) 与包含映射
\[
C(L'_1,\lambda'_1)
\longrightarrow
C^{r=r_0,s=s_0}
(L'_1,(\lambda_0,\lambda'_1))
\]
的同伦逆之复合。

\par\noindent\textbf{关于 (7.2.280.2)。}
由第 7.2.10 小节的论证，
可得滤过 $A_{n',K'}$ 代数
\[
C^{s=s_0}
([0,1]_r\times L'_1,(\lambda_0,\lambda'_1))
\]
以及滤过 $A_{n',K'}$ 同态
\begin{equation}
\operatorname{Incl}_r\colon
C(L'_1,\lambda'_1)\longrightarrow
C^{s=s_0}
([0,1]_r\times L'_1,(\lambda_0,\lambda'_1)),
\tag{7.2.282.1}\label{eq:7-2-282-1}
\end{equation}
\begin{equation}
\operatorname{Eval}_{r=r_0}\colon
C^{s=s_0}
([0,1]_r\times L'_1,(\lambda_0,\lambda'_1))
\longrightarrow
C^{r=r_0,s=s_0}
(L'_1,(\lambda_0,\lambda'_1)),
\tag{7.2.282.2}\label{eq:7-2-282-2}
\end{equation}
以及滤过 $A_{n,K}$ 同态
\begin{equation}
\mathfrak f^{\{\Phi_{r,s_0}\}_r}\colon
C(L_0,\lambda_0)\longrightarrow
C^{s=s_0}
([0,1]_r\times L'_1,(\lambda_0,\lambda'_1)),
\tag{7.2.282.3}\label{eq:7-2-282-3}
\end{equation}
使
\par\noindent\textup{(7.2.282.4)}
$\operatorname{Eval}$ 与
$\operatorname{Incl}$ 满足 (7.2.175)；
\begin{equation}
\operatorname{Eval}_{r=0}\circ
\mathfrak f^{\{\Phi_{r,s_0}\}_r}
=\mathfrak f^{\phi'\circ\psi_{0,s_0}},
\tag{7.2.282.5}\label{eq:7-2-282-5}
\end{equation}
\begin{equation}
\operatorname{Eval}_{r=1}\circ
\mathfrak f^{\{\Phi_{r,s_0}\}_r}
=\mathfrak f^{\psi_{1,s_0}\circ\phi}.
\tag{7.2.282.6}\label{eq:7-2-282-6}
\end{equation}
式 (7.2.266) 用来定义
同时满足 (7.2.282.5)、
(7.2.282.6) 的映射 (7.2.282.3)
\translatornote{原书此处把
(7.2.282.6) 误引为
(7.2.285.6)；后者不是本段列出的
端点评价式。}。

接着用同伦
$\{\Phi_{r_0,s}\}_s$
得到滤过 $A_{n',K'}$ 代数
\[
C^{r=r_0}
([0,1]_s\times L'_1,(\lambda_0,\lambda'_1))
\]
及滤过 $A_{n',K'}$ 同态
\begin{equation}
\operatorname{Incl}_s\colon
C(L'_1,\lambda'_1)\longrightarrow
C^{r=r_0}
([0,1]_s\times L'_1,(\lambda_0,\lambda'_1)),
\tag{7.2.283.1}\label{eq:7-2-283-1}
\end{equation}
\begin{equation}
\operatorname{Eval}_{s=s_0}\colon
C^{r=r_0}
([0,1]_s\times L'_1,(\lambda_0,\lambda'_1))
\longrightarrow C(L'_1,\lambda'_1),
\tag{7.2.283.2}\label{eq:7-2-283-2}
\end{equation}
以及滤过 $A_{n,K}$ 同态
\begin{equation}
\mathfrak f^{\{\Phi_{r_0,s}\}_s}\colon
C(L_0,\lambda_0)\longrightarrow
C^{r=r_0}
([0,1]_s\times L'_1,(\lambda_0,\lambda'_1)),
\tag{7.2.283.3}\label{eq:7-2-283-3}
\end{equation}
使

\par\noindent\textup{(7.2.284.1)}
$\operatorname{Eval}$ 与
$\operatorname{Incl}$ 满足 (7.2.175)；

\par\noindent\textup{(7.2.284.2)}
$\operatorname{Eval}_{s=s_0}\circ
\mathfrak f^{\{\Phi_{0,s}\}_s}
=\mathfrak f^{\phi'\circ\psi_{0,s_0}}$；

\par\noindent\textup{(7.2.284.3)}
$\operatorname{Eval}_{s=s_0}\circ
\mathfrak f^{\{\Phi_{1,s}\}_s}
=\mathfrak f^{\psi_{1,s_0}\circ\phi}$。
